\section{Cel projektu}
\label{sec:cel}

Celem projektu jest aplikacja, dzięki której kierowca ma dokumenty przewozowe
w telefonie, a nie tylko w papierowej teczce. Aplikacja ma skrócić czas od
otrzymania dokumentu do jego dostępności w formie cyfrowej, zmniejszyć ryzyko
zgubienia dokumentu i zapewnić jego poufność.

\subsection{Wymagania funkcjonalne}

\begin{enumerate}
  \item Rejestracja konta (formularz z walidacją danych) i logowanie loginem
        i hasłem; uwierzytelnianie żądań tokenem JWT.
  \item Profil kierowcy, do którego przypisane są jego dokumenty.
  \item Zarządzanie dokumentami: typ dokumentu (np. CMR, WZ), status
        (np. roboczy, przesłany, zaakceptowany), tytuł i numer, wyszukiwanie
        i filtrowanie.
  \item Zdjęcie dokumentu aparatem urządzenia mobilnego
        (interfejs \texttt{getUserMedia}), podgląd i kadrowanie przed zapisaniem
        oraz kompresja zdjęcia przed wysłaniem na serwer.
  \item Wersjonowanie dokumentów: każda nowa wersja pliku jest zachowywana,
        a zmiany metadanych i statusu trafiają do historii zmian.
\end{enumerate}

\subsection{Wymagania niefunkcjonalne}

\begin{enumerate}
  \item \textbf{Bezpieczeństwo}: pliki dokumentów są przechowywane
        w postaci zaszyfrowanej, hasła jako skróty kryptograficzne, a użytkownik
        ma dostęp wyłącznie do własnych dokumentów.
  \item \textbf{Obsługa na urządzeniach mobilnych}: interfejs responsywny,
        dostosowany do obsługi dotykowej, możliwy do zainstalowania jako PWA.
  \item \textbf{Łatwe wdrożenie}: cały system uruchamiany jednym poleceniem
        przez Docker Compose.
  \item \textbf{Niezawodność i obserwowalność}: monitorowanie stanu zdrowia
        usług, metryki, alerty i automatyczne restarty po awarii.
  \item \textbf{Jakość}: testy funkcjonalne każdej funkcji systemu,
        uzupełnione testami jednostkowymi.
\end{enumerate}
